\section{Recover signed changes from unsigned bit fields}
A DIFF tag carries three two-bit fields. Each stored field is the actual component
change plus two. Extract the field, subtract two, then add the result to the previous
component. Thus a stored zero means minus two, and stored three means plus one.
The possible changes are \(-2,-1,0,1\), independently for red, green and blue.
Alpha remains unchanged.

Component addition wraps modulo 256. From red 255, a change of plus one produces zero.
From red zero, minus one produces 255. Clamp-to-zero or saturation is not a substitute:
those operations reconstruct different component codes. Mask the updated components
to eight bits before combining the color word. Carrying an unmasked negative value
into the next field can set many unrelated bits.

For the worked tag \texttt{7f}, remove the \texttt{40} class bits and split the low
six bits into \texttt{11 11 11}. Each field means plus one. Starting with opaque black,
the result is \((1,1,1,255)\). A decoder need not recognize that this looks nearly black;
it needs the exact three additions and the retained alpha.

An encoder has the inverse problem. Compute a signed representative of each byte
change, accounting for wraparound, then test whether all three fit from minus two
through plus one. A change from zero to 255 can be represented by minus one; using
only the ordinary integer difference 255 misses that short encoding. Choosing RGB
instead still reconstructs the same pixel, but uses more bytes.
